\section*{Chapitre \ref{ch:s2:flow-market-making-at-a-bank} --- Market making sur flux dans une banque}

\begin{solution}{exo:s2:flow-market-making-at-a-bank:1}
$(1.2 - 0.83) \times \$100 = \$37$ par transaction, avant les coûts de couverture.
\end{solution}

\begin{solution}{exo:s2:flow-market-making-at-a-bank:2}
Le skew vaut $0.05 \times 4 = 0.2$ points de base contre la position longue~: un client acheteur, qui la réduit, paie $1.2 - 0.2 = 1.0$~; un client vendeur,
qui l'alourdit, paie $1.2 + 0.2 = 1.4$.
\end{solution}

\begin{solution}{exo:s2:flow-market-making-at-a-bank:3}
À partir de $0.8 = 1/(1 + e^{(1.2 - c)/0.25})$~: $c = 1.2 + 0.25 \ln(0.8/0.2) = 1.55$ points de base.
\end{solution}

\begin{solution}{exo:s2:flow-market-making-at-a-bank:4}
Un prix unique doit servir des clients aux mark-out différents et aux cotations concurrentes différentes. Un prix bas gagne les clients informés, dont les transactions perdent
de l'argent~; un prix élevé fait perdre le volume des clients captifs avec celui des informés. La tarification client par client facture aux clients captifs ce qu'ils acceptent de payer et écarte
les clients informés, ce qu'aucun chiffre unique ne peut faire.
\end{solution}

\begin{solution}{exo:s2:flow-market-making-at-a-bank:5}
Un petit skew coûte peu de capture et économise les coûts de couverture et les réévaluations de l'inventaire~: les clients négocient davantage dans le sens qui réduit la position. Un grand skew
cède de la capture sur de nombreuses transactions pour économiser des coûts de couverture déjà presque nuls~: 7.63~millions de dollars à 0.05, 7.08~millions de dollars à 0.2.
\end{solution}

\begin{solution}{exo:s2:flow-market-making-at-a-bank:6}
Un dealer peut décaler ses cours pour attirer des flux de compensation~; les clients sensibles au prix réagissent au skew, si bien que la position est compensée plus vite et à une concession plus faible
comme l'ont montré Butz et Oomen.
\end{solution}

\begin{solution}{exo:s2:flow-market-making-at-a-bank:7}
Les transactions passent à 83\,027, les réévaluations tombent à $-1.65$ millions de dollars et le \omterm{def:s2:flow-market-making-at-a-bank:franchise}{P\&L de franchise} à 4.60~millions de dollars. Les clients informés négocient désormais 3\,259 fois à un
net de $-1.72$ points de base par transaction, et même les clients modérés font perdre de l'argent au desk ($-0.07$).
\end{solution}

\begin{solution}{exo:s2:flow-market-making-at-a-bank:8}
Le volume issu des clients informés est la perte du desk~: au prix plat, leurs transactions ont dégagé $-1.18$ points de base chacune, et des cours plus serrés permettraient d'en gagner
davantage. Ce sont leurs mark-out, et non leur volume, qui doivent fixer leurs prix, et c'est le \omterm{def:s2:flow-market-making-at-a-bank:franchise}{P\&L de franchise}, non le volume, qui mesure le desk.
\end{solution}

\begin{solution}{pb:s2:flow-market-making-at-a-bank:1}
\begin{enumerate}
\item Le profit tiré du service rendu aux clients~: capture du spread, réévaluations de l'inventaire (où apparaît la sélection adverse) et coûts de couverture.
\item 60 clients non informés, 20 modérément informés (dérive de 0.8 point de base) et 10 informés (2.5), chacun avec sa propre cotation concurrente~; la dérive suit
la transaction du client, quel que soit celui qui l'exécute, de sorte que la trajectoire du mid est commune à toutes les politiques.
\item Mark-out de $-0.01$, 0.84 et 2.41 points de base et taux de touche (hit ratio) de 80~\%, 43~\% et 8~\%~; les cotations concurrentes inférées sont corrélées à 0.998 avec les
vraies.
\item Les signatures de prix, c'est-à-dire la trajectoire moyenne du prix autour des exécutions, distinguent les fournisseurs de liquidité qui internalisent de ceux qui externalisent.
\item Tarifer chaque client en fonction de son mark-out et de sa sensibilité au prix.
\item La demi-fourchette qui maximise (demi-fourchette moins mark-out) multipliée par la probabilité de gain, la cotation concurrente étant inférée à partir du taux de touche.
\item Fourchette fixe de 1.0, 1.2, 1.4, 1.6~: 5.90, 6.57, 6.63 et 6.23~millions de dollars~; tarification client par client 7.60~millions de dollars, avec le skew 7.63~millions de dollars.
\item Non informés~: capture de 1.20 à 1.39, net de 1.20 à 1.40 points de base~; modérés~: 1.20 à 1.25, net de 0.37 à 0.44~; informés~: 938 transactions à $-1.18$ à
10 transactions.
\item Un décalage de la cotation contre l'inventaire ou les axes, pour attirer des transactions qui réduisent la position.
\item À 0.05~: P\&L de 7.63~millions de dollars, inventaire moyen de 2.18 unités au lieu de 2.75, internalisation de 96.8~\% au lieu de 90.8~\%~; à 0.2~: 7.08~millions de dollars,
1.30 unités, toutes internalisées.
\item Horizons d'internalisation de plusieurs minutes sur les devises liquides et de dizaines de minutes sur les marchés émergents pour un grand dealer représentatif~;
des coûts plus faibles pour les dealers qui détiennent davantage de risque et font face à des clients sensibles au prix.
\item L'option du dealer de rejeter ou de recoter une demande de cotation après la latence~; sa conception fixe le risque d'exécution du client sans nécessairement modifier ses
coûts effectifs.
\item Avec une tarification client par client, le desk capte 1.39 point de base par transaction auprès des clients non informés et 1.25 auprès des clients modérément informés, écarte
les informés, et internalise 90.8~\% de son volume, 96.8~\% avec un skew de 0.05.
\item Que la forme de la courbe de probabilité de gain soit connue (logistique de largeur connue)~: le desk a inféré les cotations concurrentes presque exactement.
\item Il gagne les transactions des clients informés et perd sur chacune.
\item Faire varier délibérément les cotations (petits décalages aléatoires) et ajuster les taux de touche en fonction de celles-ci, par segment de clients et par taille.
\item Recoter ce client~: élargir sa fourchette ou réduire sa taille, et vérifier si son comportement a changé ou si son estimation de mark-out était bruitée.
\item Internalisation avant couverture, qui garde le risque en attendant des flux de compensation.
\item Le chapitre~15 négocie contre des flux prévisibles venus de l'extérieur~; ce chapitre est le desk qui reçoit ces flux et les tarifie.
\item L'immédiateté, tarifée selon l'information et les alternatives de chaque client.
\end{enumerate}
\end{solution}

\begin{solution}{iq:s2:flow-market-making-at-a-bank:1}
Le mouvement du mid après une transaction, mesuré en faveur du client, sur un ou plusieurs horizons. Moyenné sur les transactions d'un client, c'est ce que le desk perd
à cause de l'information de ce client~; il faut l'ajouter à la marge cible pour fixer la fourchette du client.
\end{solution}

\begin{solution}{iq:s2:flow-market-making-at-a-bank:2}
Utiliser toutes les demandes de cotation, gagnées et perdues~: le taux de touche par rapport au cours affiché. Randomiser légèrement les cotations pour que la courbe soit identifiée, ajuster un modèle de la
probabilité de gain par segment de clients, taille et conditions de marché, et corriger le fait que les cotations ont été fixées en tenant compte de l'information sur le client.
\end{solution}

\begin{solution}{iq:s2:flow-market-making-at-a-bank:3}
Une cotation biface décalée~: un cours vendeur plus serré pour réduire la position et un cours acheteur plus large, calibrés en fonction du mark-out du client, de la position et du coût de
couverture, et non un refus unilatéral, qui nuit à la franchise.
\end{solution}

\begin{solution}{iq:s2:flow-market-making-at-a-bank:4}
À partir du délai attendu avant les flux de compensation et de la volatilité sur ce délai~: la perte de la limite de risque lors d'un mouvement de stress doit tenir dans le budget de risque du desk~; il faut
la tester contre des marchés unidirectionnels, lorsque les flux de compensation n'arrivent pas.
\end{solution}

\begin{solution}{iq:s2:flow-market-making-at-a-bank:5}
S'abonner aux transactions et à un mid de référence~; pour chaque transaction, programmer les calculs de mark-out à horizons fixes~; agréger par client avec des
poids décroissants~; publier les estimations assorties de leur niveau de confiance au moteur de valorisation~; conserver les données brutes pour la recherche et l'audit.
\end{solution}

\begin{solution}{iq:s2:flow-market-making-at-a-bank:6}
Maximiser $(h - m)p(h)$~: $p + (h - m)p' = 0$, et pour la logistique $p' = -p(1-p)/w$, donc $p - (h - m)p(1-p)/w = 0$ et $h - m = w/(1 - p(h))$. La
marge au-dessus du mark-out croît avec la probabilité de gain~: les clients captifs paient davantage.
\end{solution}
